% Laget automatisk av elevark/lag_elevark.py fra sim/sprettball.html. Endringer her blir overskrevet.
\documentclass[11pt]{article}
\usepackage{elevark}
\begin{document}
\elevarkhode{Fysikk 1 · Mekanisk energi}{Sprettball}{https://lektorbo.no/sim/sprettball.html}{Simuleringen}

En ball slippes ned på et hardt gulv. Hvor høyt spretter den, og hvor mye av energien forsvinner i hvert sprett? Prøv først med simuleringen, og mål deretter på en ekte ball.

\section*{Oppgaver}

Jobb gjerne i par. Skriv ned gjetningen og begrunnelsen \emph{før} dere trykker Start.

\begin{oppgaveliste}
\oppgave \fase{gjett}{Gjett}En tennisball slippes fra 2,0~m. Hvor høyt spretter den?
\svar
\oppgave \fase{test}{Test}Les av høyden etter de tre første sprettene. Hvor stor andel av høyden beholder ballen i hvert sprett? Er andelen den samme hver gang?
\svar
\oppgave \fase{forklar}{Forklar}Hvor mange prosent av den mekaniske energien forsvinner i hvert sprett? Hvor blir det av energien?
\svar
\oppgave \fase{gjett}{Gjett}En basketball slippes fra 2,0~m. Hvor mange sprett tar det før den spretter lavere enn 10~cm? Regn ut, og test.
\svar
\oppgave \fase{test}{Mål}Gjør forsøket med en ekte ball og fyll inn tabellen under «Egne målinger». Stemmer modellen med en fast andel for ballen deres?
\svar
\oppgave \fase{utfordring}{Utfordring}Velg «Ukjent ball». Finn andelen av høyden som ballen beholder. Vis at forholdet mellom farten rett etter og rett før et sprett er $e = \sqrt{\text{andel}}$. Dette tallet kalles støttallet.
\svar
\end{oppgaveliste}

\end{document}
